% ---------------------------------------------------------------------
%  CHƯƠNG 4 — CÀI ĐẶT VÀ THỬ NGHIỆM
% ---------------------------------------------------------------------
\chapter{Cài đặt và thử nghiệm}
\label{ch:caidat}

\begin{muctieu}
Chương này mô tả môi trường triển khai, các thành phần chính kèm mã nguồn
minh hoạ, và giao diện demo của hệ thống.
\end{muctieu}

\section{Môi trường triển khai}

\begin{chuy}
\goiy{ghi rõ cấu hình phần cứng, hệ điều hành và phiên bản phần mềm để người
khác tái lập được kết quả. Kèm cả đường dẫn kho \code{git} chứa mã nguồn.}
\end{chuy}

\begin{mau}
Hệ thống được phát triển và thử nghiệm trên máy tính cá nhân với CPU Intel
Core i5, RAM 8~GB, hệ điều hành Windows 11 và Python 3.11 cùng các thư viện
đã nêu ở Bảng~\ref{tab:congcu}. Mã nguồn được quản lý bằng \code{git} và đặt
tại kho \url{https://github.com/tencuaem/ten-do-an}.
\end{mau}

\section{Các thành phần chính}

\begin{chuy}
\goiy{trình bày những đoạn mã quan trọng nhất như tiền xử lý, huấn luyện,
dự đoán. Chỉ đưa phần cốt lõi, không dán toàn bộ mã nguồn vào thân báo cáo;
mã dài để ở phụ lục. Comment trong code viết bằng tiếng Anh theo quy ước của
khoa.}
\end{chuy}

Mã nguồn~\ref{lst:train} minh hoạ bước vec-tơ hoá và huấn luyện mô hình.

\begin{lstlisting}[language=Python, caption={Huấn luyện mô hình phân lớp cảm xúc.}, label={lst:train}]
from sklearn.feature_extraction.text import TfidfVectorizer
from sklearn.linear_model import LogisticRegression
from sklearn.pipeline import Pipeline

# Build a pipeline: TF-IDF features feed a logistic regression classifier
model = Pipeline([
    ("tfidf", TfidfVectorizer(ngram_range=(1, 2), min_df=2)),
    ("clf", LogisticRegression(max_iter=1000)),
])

# Fit the model on the training split only
model.fit(X_train, y_train)
\end{lstlisting}

Mã nguồn~\ref{lst:predict} minh hoạ hàm dự đoán dùng cho giao diện web.

\begin{lstlisting}[language=Python, caption={Hàm dự đoán cảm xúc cho một câu.}, label={lst:predict}]
def predict_sentiment(text):
    # Return the predicted label together with its confidence score
    label = model.predict([text])[0]
    confidence = model.predict_proba([text]).max()
    return label, round(float(confidence), 3)
\end{lstlisting}

\section{Giao diện và demo}

\begin{chuy}
\goiy{chèn ảnh chụp màn hình giao diện chính. Đặt file ảnh vào thư mục
\code{hinh/}, rồi bỏ dấu phần trăm ở dòng \code{includegraphics} bên dưới và
xoá khung giữ chỗ.}
\end{chuy}

\begin{figure}[H]
\centering
% Bỏ dấu % ở dòng dưới và thay ten-anh.png bằng ảnh thật của em
% \includegraphics[width=0.8\textwidth]{ten-anh.png}
\fbox{\begin{minipage}[c][5cm][c]{0.8\textwidth}\centering
  \itshape\color{gray500}[Chèn ảnh chụp màn hình giao diện demo tại đây]
\end{minipage}}
\caption{Giao diện web cho phép nhập câu và xem kết quả dự đoán.}
\label{fig:demo}
\end{figure}

\section{Kiểm thử}

\begin{chuy}
\goiy{nêu cách em kiểm tra hệ thống chạy đúng: các ca kiểm thử đã chạy, dữ
liệu biên đã thử, lỗi đã phát hiện và cách khắc phục. Một bảng ca kiểm thử
ngắn cũng đủ và rất được đánh giá cao.}
\end{chuy}

\begin{mau}
Hệ thống được kiểm thử với ba nhóm ca. Nhóm một là câu ngắn rõ sắc thái,
nhóm hai là câu dài nhiều mệnh đề, nhóm ba là dữ liệu biên gồm chuỗi rỗng,
chuỗi chỉ có ký tự đặc biệt và chuỗi rất dài. Ở nhóm ba, hệ thống ban đầu
báo lỗi với chuỗi rỗng; lỗi được khắc phục bằng cách kiểm tra đầu vào trước
khi gọi mô hình.
\end{mau}
